% Vietnamese translation for OI Public Library
% Source-PDF: IOI中国国家候选队论文/1999/周咏基.pdf
\documentclass[11pt]{article}
\usepackage{oipl-vn}

\title{Bàn về nguyên lý và thiết kế thuật toán ngẫu nhiên hóa}
\author{Zhou Yongji}
\date{}

\begin{document}
\maketitle

\origtitle{On the principles and design of randomized algorithms}
\sourcepath{IOI China candidate team papers / 1999 / Zhou Yongji.pdf}

\paragraph{Từ khóa.}
Thuật toán ngẫu nhiên hóa; tính ổn định.

\section*{Tóm tắt}

Bài viết này đề xuất một loại thuật toán mới để giải các bài toán tin học:
thuật toán ngẫu nhiên hóa, đồng thời thảo luận nguyên lý và phương pháp thiết
kế của nó. Trước hết bài viết đưa ra định nghĩa thuật toán ngẫu nhiên hóa và
giải thích rằng, do ảnh hưởng của ``vận may'', ta phải phân tích tính ổn định
của thuật toán ngẫu nhiên hóa. Sau đó, bài viết chia thành ba trường hợp:
``ngẫu nhiên không ảnh hưởng tới kết quả thực thi'', ``ngẫu nhiên ảnh hưởng
tới tính đúng của kết quả thực thi'', và ``ngẫu nhiên ảnh hưởng tới chất lượng
của kết quả thực thi''. Từ các thuật toán cơ bản đến các ví dụ dùng thuật
toán ngẫu nhiên hóa để giải hiệu quả bài thi, bài viết phân tích chi tiết
nguyên lý và cách thiết kế thuật toán ngẫu nhiên hóa trong cả ba trường hợp.
Cuối cùng, bài viết tổng kết nguyên lý cơ bản và tính chất chung của thuật
toán ngẫu nhiên hóa, nêu phương pháp tổng quát để thiết kế loại thuật toán
này, đồng thời chỉ ra phạm vi áp dụng và những đặc điểm mà một thuật toán
ngẫu nhiên hóa hiệu quả cần có.

\section{Dẫn nhập}

Trong bài viết này, ta nghiên cứu một khái niệm thuật toán mới: thuật toán
ngẫu nhiên hóa. Như tên gọi, ngẫu nhiên hóa là việc dùng hàm ngẫu nhiên. Ở
đây có thể lấy hàm \texttt{RANDOM(N)} trong Borland Pascal hoặc Turbo Pascal
làm ví dụ; giá trị trả về của nó là một số nguyên trong đoạn \([0,N-1]\), và
mỗi số nguyên được trả về với xác suất bằng nhau~[1].

Một thuật toán chứa hàm ngẫu nhiên rất có thể~[2] bị chi phối bởi yếu tố bất
định. Người ta thường cho rằng một thuật toán bị yếu tố bất định chi phối chắc
chắn không phải thuật toán hiệu quả. Chính cách nghĩ ấy đã khiến thuật toán
ngẫu nhiên hóa bị xem nhẹ trong thời gian dài. Nhưng trong phần thảo luận sau,
ta sẽ thấy điều hoàn toàn ngược lại: với một số bài toán cụ thể, thuật toán
ngẫu nhiên hóa lại trở thành công cụ giải bài rất hiệu quả, đôi khi còn tốt
hơn thuật toán không ngẫu nhiên hóa thông thường.

\subsection{Định nghĩa thuật toán ngẫu nhiên hóa}

Thuật toán ngẫu nhiên hóa là thuật toán có dùng hàm ngẫu nhiên, và giá trị trả
về của hàm ngẫu nhiên trực tiếp hoặc gián tiếp ảnh hưởng tới luồng thực thi
hoặc kết quả thực thi của thuật toán.

Theo định nghĩa này, không phải mọi thuật toán dùng hàm ngẫu nhiên đều được
gọi là thuật toán ngẫu nhiên hóa. Chẳng hạn, một thuật toán có câu lệnh
\[
  i \leftarrow \texttt{RANDOM}(N),
\]
nhưng biến \(i\), ngoài việc được gán một giá trị ngẫu nhiên ở đây, không hề
xuất hiện ở nơi nào khác. Rõ ràng, nếu thuật toán này không dùng hàm ngẫu
nhiên ở chỗ khác, câu lệnh trên không thể ảnh hưởng tới luồng hoặc kết quả
thực thi, nên thuật toán đó không thể gọi là thuật toán ngẫu nhiên hóa.

Mặt khác, nếu một thuật toán là thuật toán ngẫu nhiên hóa, thì luồng thực thi
hoặc kết quả của nó sẽ chịu ảnh hưởng của hàm ngẫu nhiên được dùng trong đó.
Theo tính chất và mức độ ảnh hưởng, ta chia thành ba trường hợp:
\begin{enumerate}
  \item Ngẫu nhiên không ảnh hưởng tới kết quả thực thi. Khi ấy ngẫu nhiên tất
  yếu ảnh hưởng tới luồng thực thi, và hiệu ứng của nó thường biểu hiện qua sự
  dao động của hiệu suất thời gian.

  \item Ngẫu nhiên ảnh hưởng tới tính đúng của kết quả thực thi. Trong trường
  hợp này, bài toán gốc yêu cầu tìm một nghiệm khả thi, hoặc bài toán gốc là
  bài toán quyết định~[3]; hiệu ứng của ngẫu nhiên biểu hiện qua xác suất thu
  được nghiệm đúng khi thực thi.

  \item Ngẫu nhiên ảnh hưởng tới chất lượng của kết quả thực thi. Khi ấy hiệu
  ứng của ngẫu nhiên biểu hiện qua độ sai khác giữa kết quả thực tế và nghiệm
  tối ưu hoặc kết quả kỳ vọng trên lý thuyết.
\end{enumerate}

Trong hai trường hợp thứ hai và thứ ba, ảnh hưởng của ngẫu nhiên cũng có thể
đi kèm ảnh hưởng tới luồng thực thi. Phần sau sẽ thảo luận lần lượt theo ba
trường hợp này. Trước khi thảo luận, ta còn cần làm rõ một vấn đề.

\subsection{Ngẫu nhiên hóa và ``vận may''}

Do tình hình thực thi của thuật toán ngẫu nhiên hóa bị yếu tố bất định chi
phối, nên ngay cả cùng một thuật toán, với cùng một đầu vào trong nhiều lần
chạy, tình hình thực thi cũng sẽ khác nhau, ít nhất là hơi khác nhau. Sự khác
biệt có thể là tốc độ ra nghiệm, nghiệm đúng hay sai, nghiệm tốt hay kém,
v.v. Ví dụ, một thuật toán ngẫu nhiên hóa có thể trong hai lần chạy cho nghiệm
trước tốt hơn, nghiệm sau kém hơn. Vấn đề là trong đa số trường hợp, nhất là
trong thi đấu, với cùng một đầu vào, chương trình chỉ được chạy một lần, và
kết quả chạy được dùng để đánh giá thuật toán tốt hay kém. Như vậy ta sẽ quy
lần chạy cho nghiệm kém là do vận may không tốt, và ngược lại. Tuy nhiên, thi
đấu là so sánh thuật toán của ai hiệu quả hơn, chứ không phải ai may mắn hơn.

Một khi đã dùng hàm ngẫu nhiên, ta không thể thoát khỏi ảnh hưởng của vận may,
nên mục tiêu là hạ ảnh hưởng ấy xuống mức thấp nhất. Nói cách khác, ta phải
làm cho tình hình thực thi của thuật toán tương đối ổn định. Vì vậy trong các
phân tích thuật toán tiếp theo, ta sẽ xét hiệu năng thuật toán trên bốn mặt:
\begin{enumerate}
  \item Hiệu suất thời gian.
  \item Tính đúng của nghiệm.
  \item Mức độ tốt xấu của nghiệm, tức mức độ gần với nghiệm tối ưu.
  \item Tính ổn định, tức mức biến động trong tình hình thực thi của thuật
  toán trên cùng một đầu vào. Biến động càng nhỏ thì càng ổn định. Tính ổn
  định của thuật toán không ngẫu nhiên hóa là \(100\%\); tính ổn định của
  thuật toán ngẫu nhiên hóa thuộc khoảng \((0\%,100\%)\).
\end{enumerate}

Thông thường, miễn là cài đặt chương trình của thuật toán không vượt quá giới
hạn bộ nhớ, ta không cần cố ý nâng cao hiệu suất không gian, nên bỏ qua mục
phân tích hiệu suất không gian. ``Tính ổn định'' ở mục thứ tư có thể là độ
phức tạp thời gian trung bình của thuật toán, cũng có thể là xác suất nhận
được nghiệm đúng sau khi chạy thuật toán, và cũng có thể là xác suất nghiệm
thực tế đạt tới một mức chất lượng nào đó. ``Tính ổn định'' là một chỉ tiêu
quan trọng để đánh giá thuật toán ngẫu nhiên hóa tốt hay kém.

\section{Thuật toán ngẫu nhiên hóa có kết quả thực thi xác định}

Trong mục này, ta lấy sắp xếp nhanh và phiên bản ngẫu nhiên hóa của nó làm ví
dụ để thảo luận thuật toán ngẫu nhiên hóa có kết quả thực thi xác định. Theo
phân tích trong phần dẫn nhập, nếu kết quả thực thi của một thuật toán ngẫu
nhiên hóa là xác định, thì luồng thực thi của nó chắc chắn chịu ảnh hưởng của
ngẫu nhiên, và ảnh hưởng này thường biểu hiện qua hiệu suất thời gian. Vì vậy
trong phần dưới, ta bỏ qua phân tích về tính đúng và chất lượng của kết quả
thực thi.

\subsection{Thuật toán sắp xếp nhanh}

Sắp xếp nhanh là một phương pháp sắp xếp thường dùng. Tư tưởng cơ bản của nó
có tính đệ quy: chia dãy số cần sắp xếp thành hai phần, sao cho mỗi số ở phần
trước không lớn hơn mỗi số ở phần sau; rồi tiếp tục chia riêng từng phần cho
tới khi phần cần chia chỉ còn một số. Thuật toán có thể mô tả bằng giả mã sau.
\[
\begin{array}{ll}
\textsc{Quicksort}(A,lo,hi)\\
1 & \textbf{if } lo < hi\\
2 & \quad p \leftarrow \textsc{Partition}(A,lo,hi)\\
3 & \quad \textsc{Quicksort}(A,lo,p)\\
4 & \quad \textsc{Quicksort}(A,p+1,hi)
\end{array}
\]

Nếu \(n\) số cần sắp xếp được lưu trong mảng \(A\), gọi
\(\textsc{Quicksort}(A,1,n)\) sẽ thu được \(n\) số theo thứ tự tăng dần. Thuật
toán sắp xếp nhanh trên phụ thuộc vào thủ tục chia
\(\textsc{Partition}(A,lo,hi)\). Trong thời gian \(\Theta(n)\), thủ tục này
chia \(A[lo..hi]\) thành hai phần: các phần tử không lớn hơn \(x=A[lo]\), và
các phần tử không nhỏ hơn \(x=A[lo]\). Hai phần này lần lượt được đặt trong
\(A[lo..p]\) và \(A[p+1..hi]\). Trong thủ tục
\(\textsc{Quicksort}(A,lo,hi)\), ta gọi đệ quy \(\textsc{Quicksort}()\) để
tiếp tục chia \(A[lo..p]\) và \(A[p+1..hi]\).

Có thể chứng minh~[4] rằng trong trường hợp xấu nhất, chẳng hạn mỗi lần chia
đều làm \(p=lo\), độ phức tạp thời gian của sắp xếp nhanh là
\(\Theta(n^2)\); trong trường hợp tốt nhất, độ phức tạp thời gian là
\(\Theta(n\log_2 n)\). Nếu giả sử mọi hoán vị của đầu vào xuất hiện với xác
suất bằng nhau, điều thực tế thường không đúng, thì độ phức tạp thời gian
trung bình của thuật toán là \(O(n\log_2 n)\).

\subsection{Sắp xếp nhanh ngẫu nhiên hóa}

Qua phân tích, ta thấy sắp xếp nhanh là phương pháp sắp xếp rất hiệu quả, với
độ phức tạp thời gian trung bình \(O(n\log_2 n)\). Nhưng trong trường hợp xấu
nhất, độ phức tạp thời gian của nó là \(\Theta(n^2)\), nên khi \(n\) lớn thì
chạy rất chậm; xem bảng đối chiếu hiệu năng ở cuối mục này. Thật ra, nếu theo
giả thiết ở trên rằng mọi hoán vị đầu vào xuất hiện với xác suất bằng nhau,
xác suất gặp trường hợp xấu nhất nhỏ tới mức chỉ là \(\Theta(1/n!)\), và hằng
số ẩn trong \(\Theta()\) rất nhỏ. Nhìn theo cách đó, sắp xếp nhanh vẫn rất có
giá trị.

Nhưng thực tế thường không phù hợp với giả thiết này. Với một bài toán nào
đó, phần lớn đầu vào ta gặp có thể là trường hợp xấu nhất hoặc gần xấu nhất.
Một cách xử lý là không dùng \(x=A[lo]\) để chia \(A[lo..hi]\), mà dùng
\(x=A[hi]\), hoặc \(x=A[(lo+hi)\operatorname{div}2]\), hoặc một phần tử khác
trong \(A[lo..hi]\) để chia \(A[lo..hi]\), tùy tình huống cụ thể. Nhưng cách
này chưa giải quyết vấn đề, vì ta có thể gặp loại đầu vào sau: có ba lớp đầu
vào, mỗi lớp xuất hiện với xác suất \(1/3\), và ba lớp ấy lần lượt là trường
hợp xấu nhất đối với \(x=A[lo]\), \(x=A[hi]\), và
\(x=A[(lo+hi)\operatorname{div}2]\). Khi đó sắp xếp nhanh sẽ rất kém hiệu quả.

Sau khi ngẫu nhiên hóa sắp xếp nhanh, ta có thể khắc phục dạng vấn đề này.
Tư tưởng của sắp xếp nhanh ngẫu nhiên hóa là: trong mỗi lần chia, chọn ngẫu
nhiên một số trong \(A[lo..hi]\) làm \(x\) để chia \(A[lo..hi]\). Chỉ cần sửa
nhẹ thuật toán gốc. Ta chỉ thêm hàm \(\textsc{Partition\_R}\), hàm này gọi
thủ tục \(\textsc{Partition}()\) ban đầu. Phần sửa đổi trong
\(\textsc{Quicksort\_R}()\) so với \(\textsc{Quicksort}\) chính là việc gọi
hàm chia ngẫu nhiên.
\[
\begin{array}{ll}
\textsc{Partition\_R}(A,lo,hi)\\
1 & r \leftarrow \texttt{RANDOM}(hi-lo+1)+lo\\
2 & \text{đổi chỗ } A[lo] \text{ và } A[r]\\
3 & \textbf{return } \textsc{Partition}(A,lo,hi)
\end{array}
\]
\[
\begin{array}{ll}
\textsc{Quicksort\_R}(A,lo,hi)\\
1 & \textbf{if } lo < hi\\
2 & \quad p \leftarrow \textsc{Partition\_R}(A,lo,hi)\\
3 & \quad \textsc{Quicksort\_R}(A,lo,p)\\
4 & \quad \textsc{Quicksort\_R}(A,p+1,hi)
\end{array}
\]

\subsection{Phân tích sắp xếp nhanh ngẫu nhiên hóa}

Ngẫu nhiên hóa không thay đổi thủ tục chia của sắp xếp nhanh ban đầu, nên hiệu
suất thời gian của sắp xếp nhanh ngẫu nhiên hóa vẫn phụ thuộc vào vị trí của
số được chọn để chia trong mảng sau khi sắp xếp. Độ phức tạp thời gian xấu
nhất, trung bình và tốt nhất vẫn lần lượt là \(\Theta(n^2)\),
\(O(n\log_2 n)\), và \(\Theta(n\log_2 n)\); chỉ có điều trường hợp xấu nhất và
tốt nhất đã thay đổi. Chúng không còn do đầu vào quyết định nữa, mà do hàm
ngẫu nhiên quyết định. Nói cách khác, ta không thể đưa ra một đầu vào xấu
nhất để làm lần thực thi rơi vào trường hợp xấu nhất, trừ khi vận may không
tốt.

Như đã nêu trong phần dẫn nhập, bây giờ ta phân tích tính ổn định của sắp xếp
nhanh ngẫu nhiên hóa. Theo giả thiết mọi hoán vị xuất hiện với xác suất bằng
nhau, giả thiết này không nhất thiết đúng, sắp xếp nhanh gặp trường hợp xấu
nhất với khả năng \(\Theta(1/n!)\). Giả sử \(\texttt{RANDOM}(n)\) sinh ra
\(n\) số với xác suất bằng nhau, giả thiết này hầu như chắc chắn đúng, thì sắp
xếp nhanh ngẫu nhiên hóa cũng gặp trường hợp xấu nhất với khả năng
\(\Theta(1/n!)\). Nếu \(n\) đủ lớn, ta có xác suất hơn \(99\%\) để ``gặp may''.
Nói cách khác, thuật toán sắp xếp nhanh ngẫu nhiên hóa có tính ổn định rất
cao.

Bảng sau đối chiếu hiệu năng của sắp xếp nhanh ban đầu và sắp xếp nhanh sau
khi ngẫu nhiên hóa.
\[
\begin{array}{lll}
\hline
\text{Mục phân tích} & \text{Thuật toán gốc} & \text{Thuật toán ngẫu nhiên hóa}\\
\hline
\text{Lý thuyết: xấu nhất} & \Theta(n^2) & \Theta(n^2)\\
\text{Lý thuyết: tốt nhất} & \Theta(n\log_2 n) & \Theta(n\log_2 n)\\
\text{Lý thuyết: trung bình} & O(n\log_2 n) & O(n\log_2 n)\\
\text{Lý thuyết: ổn định} & \Theta(1) & \Theta(1-1/n!)\\
\text{Thực chạy: đầu vào ngẫu nhiên }(n=30000) & 0.22\text{s} & 0.27\text{s}\\
\text{Thực chạy: đầu vào xấu nhất }(n=30000) & 66\text{s} & 0.22\text{s}\\
\text{Thực chạy: ổn định }(n\text{ đủ lớn}) & 100\% & >99\%\\
\text{Nguyên nhân xấu nhất} & \text{đầu vào xấu nhất} & \text{giá trị ngẫu nhiên không tốt}\\
\text{Phụ thuộc thời gian vào đầu vào} & \text{phụ thuộc hoàn toàn} & \text{không phụ thuộc}\\
\hline
\end{array}
\]

Vài chú thích cho bảng trên:
\begin{enumerate}
  \item Môi trường chạy chương trình là Pentium 100MHz, biên dịch bằng BP7.0.
  \item Thời gian chạy chương trình tương ứng của thuật toán ngẫu nhiên hóa là
  trung bình của \(1000\) lần chạy.
  \item Khi kiểm tra tính ổn định của thuật toán ngẫu nhiên hóa, chương trình
  tương ứng được chạy \(1000\) lần trên từng đầu vào khác nhau.
  \item Mã chương trình xem trong \texttt{QSORT.PAS}.
\end{enumerate}

\subsection{Tiểu kết}

Từ phân tích trên có thể thấy nguyên lý của thuật toán ngẫu nhiên hóa có kết
quả thực thi xác định là: dùng hàm ngẫu nhiên để triệt tiêu toàn bộ hoặc một
phần tác dụng của đầu vào xấu nhất, khiến hiệu suất thời gian của thuật toán
không hoàn toàn phụ thuộc vào đầu vào tốt hay xấu.

Điều khiển đầu vào một cách thích hợp để làm kết quả thực thi tương đối ổn
định là phương pháp thường dùng để thiết kế loại thuật toán ngẫu nhiên hóa
này. Ví dụ, trong sắp xếp nhanh ngẫu nhiên hóa, mỗi lần ta chọn ngẫu nhiên
\(x\) để chia \(A[lo..hi]\). Hiệu ứng của phương pháp này tương đương với
việc xáo ngẫu nhiên các số trong \(A\) trước khi sắp xếp. Một ví dụ khác là
khi xây dựng cây nhị phân tìm kiếm, có thể xáo ngẫu nhiên thứ tự các khóa cần
chèn trước, rồi lần lượt chèn vào cây, để thu được cây tìm kiếm nhị phân cân
bằng hơn và nâng cao hiệu quả tìm kiếm khóa về sau.

\section{Thuật toán ngẫu nhiên hóa có kết quả thực thi có thể lệch khỏi nghiệm đúng}

Trong mục này ta thảo luận loại thuật toán ngẫu nhiên hóa thứ hai. Loại thuật
toán này thậm chí có thể xuất kết quả sai, nhưng vẫn rất hiệu quả. Ta lấy
thuật toán kiểm tra số nguyên tố làm ví dụ.

\subsection{Thuật toán kiểm tra nguyên tố thô sơ}

Với \(n\) nhỏ, ta có thể dùng ``phương pháp sàng'' để quyết định \(n\) có là
số nguyên tố hay không. Với \(n\) lớn hơn một chút, ta có thể tìm trước mọi số
nguyên tố trong đoạn \([2,\lfloor\sqrt n\rfloor]\), rồi dùng các số nguyên tố
ấy thử chia \(n\). Hai phương pháp này đều phải dựa vào mảng lớn; nếu \(n\)
đủ lớn thì không còn thích hợp. Khi ấy ta chỉ có thể dùng
\(2,3,\ldots,\lfloor\sqrt n\rfloor\) để thử chia \(n\). Một khi chia hết, \(n\)
chắc chắn là hợp số; ngược lại là số nguyên tố. Thuật toán được mô tả như sau:
\[
\begin{array}{ll}
\textsc{IsPrime\_Naive}(n)\\
1 & \textbf{for } a \leftarrow 2 \textbf{ to } \lfloor\sqrt n\rfloor\\
2 & \quad \textbf{if } a\mid n\\
3 & \quad\quad \textbf{return } \textsc{False}\\
4 & \textbf{return } \textsc{True}
\end{array}
\]

Khi cài đặt, ta có thể kiểm tra trước \(n\) có chẵn hay không, rồi dùng
\(3,5,7,9,\ldots\) thử chia \(n\), để tăng tốc chương trình. Dù vậy, khi gặp
số nguyên tố \(n\) khá lớn, thuật toán này vẫn tỏ ra rất chậm. Độ phức tạp
thời gian xấu nhất của nó là \(\Theta(n^{1/2})\).

\subsection{Thuật toán kiểm tra nguyên tố ngẫu nhiên hóa}

Đổi góc nhìn, từ định lý Fermat ta biết: nếu \(n\) là số nguyên tố và \(a\)
không chia hết \(n\), thì nhất định có
\[
  a^{n-1}\equiv 1 \pmod n.
\]
Ta viết lại thành: nếu \(n\) là số nguyên tố, thì với
\(a=1,2,\ldots,n-1\), có \(a^{n-1}\equiv1\pmod n\). Vì vậy, nếu tồn tại số
nguyên \(a\in[1,n-1]\) sao cho \(a^{n-1}\not\equiv1\pmod n\), thì \(n\) chắc
chắn là hợp số. Xét thuật toán sau:
\[
\begin{array}{ll}
\textsc{IsPrime\_R}(n,s)\\
1 & \textbf{for } i \leftarrow 1 \textbf{ to } s\\
2 & \quad a \leftarrow \texttt{RANDOM}(n-1)+1\\
3 & \quad \textbf{if } a^{n-1}\not\equiv 1 \pmod n\\
4 & \quad\quad \textbf{return } \textsc{False}\\
5 & \textbf{return } \textsc{True}
\end{array}
\]

Thuật toán ngẫu nhiên chọn \(s\) giá trị \(a\), kiểm tra
\(a^{n-1}\equiv1\pmod n\) có đúng hay không. Nếu tìm thấy một \(a\) làm công
thức không đúng, thuật toán chắc chắn kết luận ``\(n\) là hợp số''; nếu mọi
\(a\) được chọn đều làm \(a^{n-1}\equiv1\pmod n\) đúng, thuật toán đưa ra giả
thiết ``\(n\) là số nguyên tố''.

Thuật toán này chỉ sinh một loại lỗi: khi \(s\) giá trị \(a\) được chọn đều
thỏa \(a^{n-1}\equiv1\pmod n\) nhưng \(n\) thật ra là hợp số, thuật toán sẽ
cho rằng chưa đủ bằng chứng \(n\) là hợp số và phán \(n\) là số nguyên tố.

Dòng 3 của giả mã cần tính \(a^{n-1}\bmod n\). Ta chỉ cần
\(\lfloor\log_2(n-1)\rfloor\) vòng lặp để tính giá trị này. Đặt
\[
  n-1=b_k2^k+b_{k-1}2^{k-1}+\cdots+b_0 2^0,
\]
trong đó \(b_kb_{k-1}\ldots b_0\) là biểu diễn nhị phân của \(n-1\), nên
\(k=\lfloor\log_2(n-1)\rfloor\). Khi đó
\[
  a^{n-1}\bmod n
  = a^{(\cdots(((b_k)\cdot2+b_{k-1})\cdot2+b_{k-2})\cdot2+\cdots+b_1)\cdot2+b_0}
    \bmod n.
\]
Ký hiệu \(\langle a\rangle=a\bmod n\), biểu thức trên có thể viết lại theo
dạng bình phương lặp, và dùng \(k\) vòng lặp là tính được giá trị đó. Vì vậy
độ phức tạp thời gian xấu nhất của thuật toán kiểm tra nguyên tố ngẫu nhiên
hóa là \(\Theta(s\log_2 n)\). Nếu xem \(s\) là hằng số, độ phức tạp thời gian
là \(\Theta(\log_2 n)\).

\subsection{So sánh hai thuật toán}

Tính ổn định của thuật toán kiểm tra nguyên tố ngẫu nhiên hóa phụ thuộc vào
số lượng giá trị \(a\) trong \([1,n-1]\) thỏa \(a^{n-1}\equiv1\pmod n\), với
\(n\) là hợp số. Ký hiệu số lượng đó là \(H(n)\), thì tính ổn định khi thuật
toán phán định \(n\) là \((1-H(n)/n)^s\). Trên thực tế, với đa số \(n\),
\(H(n)/n\) có thể đạt \(98\%\); với hầu hết mọi \(n\), \(H(n)/n\) không nhỏ
hơn \(50\%\); trong phạm vi \(10000\), số \(n\) có \(H(n)/n<50\%\) không quá
\(10\). Ta có lý do để tin rằng chỉ cần chọn \(s\) thích hợp là có thể làm
thuật toán có tính ổn định rất cao. Chẳng hạn lấy \(s=50\), thuật toán có
tính ổn định ít nhất \(1-2^{-50}>99\%\) với hầu hết mọi \(n\). Ngay cả khi
lấy \(s\) nhỏ hơn, như \(s=5\), thuật toán cũng chưa chắc thu được kết quả
sai.

Hiệu suất thời gian của thuật toán kiểm tra nguyên tố ngẫu nhiên hóa cao hơn
nhiều so với thuật toán thô sơ; điều này có thể thấy từ độ phức tạp thời gian
của chúng. Trong ứng dụng thực tế, lấy \(s=50\), \(n=761838257287\), là một số
nguyên tố, và chạy chương trình tương ứng của hai thuật toán \(100\) lần, thì
thuật toán trước cần \(20\) giây, thuật toán sau cần \(102\) giây~[5]. Thật ra
kiểm tra nguyên tố thường chỉ được gọi như một chương trình con, và số lần
dùng thực tế có thể còn hơn \(100\) lần. Có thể thấy chênh lệch giữa hai thuật
toán là rõ ràng.

\subsection{Tiểu kết}

Thuật toán ngẫu nhiên hóa có kết quả thực thi có thể lệch khỏi nghiệm đúng
dựa trên nguyên lý sau: mức độ xấp xỉ của một thuật toán gần đúng chịu ảnh
hưởng của một tham số nào đó. Khi tham số ấy lấy một giá trị đặc biệt, thuật
toán có thể thu được nghiệm đúng. Vì vậy, chỉ cần chạy thuật toán nhiều lần,
mỗi lần cho tham số lấy giá trị ngẫu nhiên trong một phạm vi chỉ định, thuật
toán có khả năng thu được kết quả đúng.

Thông thường, ta có thể thiết kế loại thuật toán ngẫu nhiên hóa này như sau:
trước hết chọn một thuật toán gần đúng, ở đây là thuật toán phán định dựa trên
định lý Fermat; sau đó thêm điều khiển ngẫu nhiên vào thuật toán, ở đây là
\(a\); cuối cùng thêm vòng lặp ngoài để chạy thuật toán gần đúng nhiều lần,
ở đây là \(s\), qua đó tạo thành thuật toán ngẫu nhiên hóa. Có thể dùng việc
lặp lại thuật toán gần đúng để điều khiển tính ổn định của thuật toán, khiến
nó đạt mức kỳ vọng.

\section{Thuật toán ngẫu nhiên hóa có chất lượng kết quả biến đổi}

Cuối cùng, ta thảo luận loại thuật toán ngẫu nhiên hóa thứ ba. Loại thuật
toán này phần lớn nhằm vào các bài toán chỉ yêu cầu nghiệm tương đối tốt, ví
dụ bài ``Tính toán song song'' trong NOI 1998.

\subsection{Một số thuật toán tham lam}

Nội dung đại ý của bài ``Tính toán song song'' là: nhập một biểu thức số học
bốn phép tính gồm \(+\), \(-\), \(\times\), \(\div\), \((\), \()\), và biến
viết bằng chữ cái in hoa; xuất một đoạn lệnh để điều khiển một máy tính song
song có hai bộ tính toán tính đúng giá trị biểu thức trong thời gian càng ngắn
càng tốt. Điểm của chương trình sẽ phụ thuộc vào mức độ tốt xấu khi so sánh
nghiệm tốt nhất mà chương trình tìm được với đáp án chuẩn, đáp án chuẩn chưa
chắc là tối ưu.

Nếu dùng thuật toán tìm kiếm để giải bài này, mỗi lần mở rộng cây tìm kiếm
phải xác định toán hạng hiện tại, toán tử và bộ tính toán, nên lượng tìm kiếm
lớn kinh khủng. Vì bài toán không yêu cầu tìm nghiệm tối ưu, ta có thể dùng
thuật toán tham lam để tìm nghiệm tương đối tốt. Tiêu chuẩn tham lam có nhiều
cách chọn, chẳng hạn mỗi lần ưu tiên chọn toán tử tốn nhiều thời gian, hoặc
mỗi lần chọn bộ tính toán rảnh sớm nhất. Hiện tại ta chưa tìm được một tiêu
chuẩn tham lam áp dụng phổ biến, và khả năng tìm được tiêu chuẩn như vậy không
lớn. Mặt khác, mỗi tiêu chuẩn hiện có chỉ có thể thu được kết quả tốt hơn trên
một số đầu vào ở một mức độ nhất định; ta hoàn toàn có thể thiết kế riêng các
đầu vào không phù hợp với từng tiêu chuẩn tham lam, khiến nó xuất kết quả
không lý tưởng. Bị giới hạn bởi các điều trên, thuật toán tham lam thuần túy
không thể giải hiệu quả bài ``Tính toán song song''.

Một cách xử lý là: do thuật toán tham lam có hiệu suất thời gian rất cao, ta
có thể thử tất cả các tiêu chuẩn tham lam trong cùng một chương trình, nhưng
điều đó chắc chắn làm tăng rất mạnh ``độ phức tạp lập trình''. Thuật toán tham
lam ngẫu nhiên hóa dưới đây đưa ra một cách giải quyết tốt hơn.

\subsection{Thuật toán tham lam ngẫu nhiên hóa}

Tư tưởng cơ bản của thuật toán tham lam ngẫu nhiên hóa là: đặt mức độ tham lam
\(\textit{rate}\%\), với \(\textit{rate}\in[0,100]\); chọn một tiêu chuẩn tham
lam tương đối tốt làm nền; mỗi lần cần tìm nghiệm tối ưu cục bộ, đổi thành tìm
một nghiệm cục bộ tương đối tốt có mức độ tham lam theo tiêu chuẩn ấy không
nhỏ hơn \(\textit{rate}\%\). Sửa đổi này có thể mô tả bằng giả mã sau:
\[
\begin{array}{ll}
\textsc{PartOptimize}(\textit{rate})\\
1 & \textbf{for } A \leftarrow \text{nghiệm tối ưu cục bộ}
    \textbf{ to } \text{nghiệm tệ nhất cục bộ}\\
2 & \quad \textbf{if } \texttt{RANDOM}(101) \le \textit{rate}\\
3 & \quad\quad \textbf{return } A\\
4 & \textbf{return } \text{nghiệm tệ nhất cục bộ}
\end{array}
\]

Với bất kỳ tiêu chuẩn tham lam hiện có nào, đều tồn tại đầu vào không phù hợp
với nó. Nói cách khác, tồn tại đầu vào khiến thuật toán, trong trường hợp
không phải lần nào cũng chọn nghiệm tối ưu cục bộ, thu được nghiệm tốt hơn so
với việc lần nào cũng chọn nghiệm tối ưu cục bộ. Thuật toán tham lam ngẫu
nhiên hóa trên có thể bao phủ tình huống này. Đồng thời, khi
\(\textit{rate}=100\), kết quả của thuật toán tham lam ngẫu nhiên hóa chính là
kết quả của thuật toán tham lam chưa sửa đổi ban đầu, nên thuật toán ngẫu
nhiên hóa trên ít nhất không kém thuật toán không ngẫu nhiên hóa.

Tư tưởng trên tương đối đơn giản, nên không khó cài đặt. Hơn nữa, các thuật
toán tham lam đều có hiệu suất thời gian cao, nên thời gian tiêu tốn cho nhiều
lần tham lam cũng không dài. Mã chương trình xem trong \texttt{PARALLEL.PAS}.
Trong cài đặt, ta còn thêm vài tối ưu khác, chẳng hạn các hạng tử lặp chỉ tính
một lần.

\subsection{Hiệu năng của thuật toán tham lam ngẫu nhiên hóa}

Do bài toán không có nhiều ràng buộc trên nghiệm, chủng loại và số lượng nghiệm
có thể rất nhiều, nên khó phân tích hiệu năng của thuật toán tham lam ngẫu
nhiên hóa trên lý thuyết. Tuy nhiên, ta có thể dùng dữ liệu kiểm thử của kỳ
thi khi đó để kiểm tra thuật toán. Cách làm này vẫn có sức thuyết phục nhất
định. Bảng sau đối chiếu kết quả chạy của các chương trình tương ứng với các
thuật toán tham lam và đáp án chuẩn.
\[
\begin{array}{lllll}
\hline
\text{Tệp vào} & \text{Tham lam} & \text{Tham lam ngẫu nhiên hóa}
  & \text{Chuẩn} & \text{Chênh lệch}\\
\hline
\texttt{INPUT.001} & 14 & 14 & 14 & 0\\
\texttt{INPUT.002} & 50 & 50 & 50 & 0\\
\texttt{INPUT.003} & 55 & 55 & 55 & 0\\
\texttt{INPUT.004} & 22 & 21 & 21 & 0\\
\texttt{INPUT.005} & 53 & 47(95\%)\;46(5\%) & 46 & +1\\
\texttt{INPUT.006} & 42 & 23 & 22 & +1\\
\texttt{INPUT.007} & 130 & 100 & 100 & 0\\
\texttt{INPUT.008} & 39 & 33 & 33 & 0\\
\texttt{INPUT.009} & 3800 & 3700 & 3700 & 0\\
\texttt{INPUT.010} & 230 & 210 & 210 & 0\\
\texttt{INPUT.011} & 10 & 12 & 10 & +2\\
\texttt{INPUT.012} & 6300 & 6300 & 6300 & 0\\
\texttt{INPUT.013} & 234 & 234(99\%)\;235(1\%) & 234 & 0\\
\texttt{INPUT.014} & 3597 & \begin{array}{l}3474(42\%)\;3486(30\%)\\3462(24\%)\;3459(4\%)\end{array} & 3498 & -24\\
\texttt{INPUT.015} & 412 & \begin{array}{l}353(52\%)\;354(22\%)\;356(16\%)\\357(6\%)\;359(3\%)\;360(1\%)\end{array} & 370 & -17\\
\texttt{INPUT.016} & 1250 & 1150(98\%)\;1250(2\%) & 1150 & 0\\
\texttt{INPUT.017} & 580 & 559(96\%)\;580(4\%) & 559 & 0\\
\texttt{INPUT.018} & 3108 & 3108 & 3108 & 0\\
\texttt{INPUT.019} & 0 & 0 & 0 & 0\\
\texttt{INPUT.020} & 192 & 192 & 171 & +21\\
\hline
\end{array}
\]

Vài chú thích cho bảng trên:
\begin{enumerate}
  \item Môi trường chạy chương trình là Pentium 100MHz, BP7.0.
  \item Các số trong cột ``thuật toán tham lam thông thường'' là nghiệm tương
  đối tốt thu được bởi thuật toán dưới nhiều tiêu chuẩn tham lam khác nhau.
  \item Chương trình tương ứng của thuật toán ngẫu nhiên hóa được chạy \(100\)
  lần để xác định kết quả chạy. Trong cột ``thuật toán tham lam ngẫu nhiên
  hóa'', xác suất thu được kết quả đó được ghi trong ngoặc cạnh số. Nếu không
  có ngoặc, nghĩa là cả \(100\) lần chạy đều cho kết quả này.
  \item Chênh lệch giữa thuật toán tham lam ngẫu nhiên hóa và đáp án chuẩn
  bằng nghiệm tương đối tốt do thuật toán tham lam ngẫu nhiên hóa thu được trừ
  đáp án chuẩn.
  \item Mã chương trình xem trong \texttt{PARALLEL.PAS}.
\end{enumerate}

Từ bảng trên, ta thấy thuật toán tham lam ngẫu nhiên hóa có thể thu được kết
quả tốt ngang đáp án chuẩn trên phần lớn đầu vào, thậm chí trên một vài đầu
vào, như \texttt{INPUT.014} và \texttt{INPUT.015}, còn thu được kết quả tốt
hơn đáp án chuẩn. Đồng thời, thuật toán này cũng có tính ổn định khá cao. Chú
ý rằng tính ổn định ở đây là xác suất thu được nghiệm trong một phạm vi nào
đó, chẳng hạn xác suất thu được kết quả tốt ngang đáp án chuẩn, hoặc xác suất
thu được kết quả hơi tốt hơn đáp án chuẩn. Ngoài ra, hiệu suất thời gian của
nó tất nhiên sẽ không kém. Có thể thấy, thuật toán tham lam ngẫu nhiên hóa ở
đây là một thuật toán hiệu quả để giải bài ``Tính toán song song''.

\subsection{Tiểu kết}

Nguyên lý của thuật toán ngẫu nhiên hóa có chất lượng kết quả biến đổi về cơ
bản giống mục trước: mức độ xấp xỉ của một thuật toán gần đúng chịu ảnh hưởng
của một tham số nào đó; khi tham số thay đổi, kết quả thực thi của thuật toán
sẽ có biến động tốt xấu. Chỉ cần chạy thuật toán nhiều lần, mỗi lần cho tham
số lấy giá trị ngẫu nhiên trong phạm vi chỉ định, và cập nhật kịp thời nghiệm
tốt nhất hiện tại sau mỗi lần chạy, nghiệm hiện tại sẽ không ngừng tiến gần
nghiệm tối ưu lý thuyết.

Thiết kế thuật toán như vậy thường lấy thuật toán tham lam làm nền và cải tạo
theo kiểu hàm \(\textsc{PartOptimize}()\). Thật ra, \(\textsc{PartOptimize}()\)
hoàn toàn có thể dùng như một khung thuật toán. Thuật toán bên ngoài thủ tục
đó có thể mô tả bằng giả mã sau:
\[
\begin{array}{ll}
\textsc{Greedy\_R}()\\
1 & bA \leftarrow \text{nghiệm tệ nhất}\\
2 & \textbf{for } \textit{rate} \leftarrow 0 \textbf{ to } 100\\
3 & \quad A \leftarrow \emptyset\\
4 & \quad \textbf{while } \text{chưa thu được nghiệm toàn cục}\\
5 & \quad\quad a \leftarrow \textsc{PartOptimize}(\textit{rate})\\
6 & \quad\quad A \leftarrow A\cup\{a\}\\
7 & \quad \textbf{if } A \text{ tốt hơn } bA\\
8 & \quad\quad bA \leftarrow A\\
9 & \textbf{return } bA
\end{array}
\]

Từ ví dụ bài ``Tính toán song song'', ta thấy khi giải những bài chỉ cần tìm
nghiệm tương đối tốt, như các bài ``Rết có hại'', ``Ghi địa danh trên bản
đồ'' trong IOI 1997, hoặc ``Đặt biển trạm'' trong CTSC 1998, thuật toán ngẫu
nhiên hóa có chất lượng kết quả biến đổi vẫn là một loại thuật toán hiệu quả.

\section{Tổng kết}

Sau các phân tích trên về ba trường hợp thuật toán ngẫu nhiên hóa, ta tổng
kết như sau.

\subsection{Nguyên lý và thiết kế thuật toán ngẫu nhiên hóa}

Nguyên lý cơ bản của thuật toán ngẫu nhiên hóa là: khi trong một quyết định có
nhiều lựa chọn nhưng không thể xác định lựa chọn nào là tốt, hoặc phải trả giá
lớn mới xác định được lựa chọn tốt, nếu đa số lựa chọn là tốt thì chọn ngẫu
nhiên một lựa chọn thường là một chiến lược hiệu quả. Thông thường, khi một
thuật toán cần đưa ra nhiều quyết định, hoặc cần chạy một thuật toán nhiều
lần, điểm này biểu hiện đặc biệt rõ.

Nguyên lý này khiến thuật toán ngẫu nhiên hóa có một tính chất chung: không có
một đầu vào đặc biệt nào có thể làm thuật toán rơi vào trường hợp xấu nhất khi
thực thi. Trường hợp xấu nhất có thể biểu hiện trong luồng thực thi, chủ yếu
là hiệu suất thời gian, cũng có thể biểu hiện trong kết quả thực thi.

Theo nguyên lý này, khi thiết kế thuật toán ngẫu nhiên hóa, ta thường lấy một
thuật toán nào đó làm mẫu, thêm yếu tố ngẫu nhiên, để thuật toán chọn ngẫu
nhiên khi khó đưa ra quyết định. Khi thiết kế thuật toán có kết quả chịu ảnh
hưởng của ngẫu nhiên, ta còn có thể chạy thuật toán nhiều lần, để thuật toán
không ngừng tiến gần nghiệm đúng hoặc nghiệm tối ưu~[6].

Trên đây chỉ là cách cơ bản để thiết kế thuật toán ngẫu nhiên hóa. Trong thực
tế, thiết kế thuật toán ngẫu nhiên hóa không có công thức có thể rập khuôn.
Ta phải phân tích từng bài toán cụ thể, nghiên cứu sâu, và thiết kế được thuật
toán ngẫu nhiên hóa tốt. Đồng thời trong thiết kế, ta còn cần đặc biệt chú ý
một vấn đề sau.

\subsection{Phạm vi áp dụng và tính hiệu quả của thuật toán ngẫu nhiên hóa}

Cuối cùng ta xét một vấn đề rất quan trọng đối với thuật toán ngẫu nhiên hóa:
phạm vi áp dụng và tính hiệu quả của nó.

Ngoài bản thân mô tả bài toán, mọi yêu cầu của một bài toán đối với thuật toán
còn gồm các yêu cầu như giới hạn bộ nhớ, giới hạn thời gian, giới hạn tính ổn
định, v.v. Nếu một bài toán không bắt buộc thuật toán phải ổn định \(100\%\),
tức với cùng một đầu vào, không nhất thiết lần chạy nào cũng giống nhau, dĩ
nhiên sự bất ổn định này không được quá lớn; đồng thời, để bù lại, bài toán
lại yêu cầu thuật toán có hiệu năng tốt ở một vài mặt khác, chẳng hạn ra
nghiệm nhanh, mà các hiệu năng này thuật toán không ngẫu nhiên hóa thông
thường không đạt được, thì khi đó thuật toán ngẫu nhiên hóa có thể là một
trong các ứng viên giải bài hiệu quả. Nếu không, thuật toán ngẫu nhiên hóa
không thích hợp.

Khi một thuật toán ngẫu nhiên hóa thích hợp để giải một bài toán, có tính ổn
định tương đối cao, đồng thời có biểu hiện nổi bật ở một vài mặt khác, chẳng
hạn tốc độ nhanh, mã ngắn, v.v., và có thể làm tốt hơn thuật toán không ngẫu
nhiên hóa thông thường, thì thuật toán ngẫu nhiên hóa ấy là một thuật toán
thật sự hiệu quả.

\section*{Phụ lục}

\paragraph{[1]}
Nói nghiêm ngặt, các số do \(\texttt{RANDOM}(N)\) trả về là số hỗn độn, không
phải số ngẫu nhiên thật sự.

\paragraph{[2]}
Ở đây không thể nói ``nhất định''. Mục định nghĩa thuật toán ngẫu nhiên hóa đã
giải thích điểm này.

\paragraph{[3]}
Bài toán quyết định là bài toán yêu cầu phán định dữ liệu đầu vào, hoặc nếu
bài toán không có đầu vào thì xem các điều kiện đã biết trong bài toán là đầu
vào, có thỏa một điều kiện đặc biệt nào đó hay không.

\paragraph{[4]}
Quá trình chứng minh như sau.

Ký hiệu \(T(n)\) là độ phức tạp thời gian của
\(\textsc{Quicksort}(A,lo,hi=lo+n-1)\). Rõ ràng \(T(n)\) phụ thuộc vào vị trí
của \(x=A[lo]\), được dùng khi chia, trong \(A[lo..hi]\) sau khi sắp xếp. Nếu
giá trị trả về của \(\textsc{Partition}(A,lo,hi)\) là \(p\), thì
\[
  T(n)=T(p-lo+1)+T(hi-p)+\Theta(n),
\]
trong đó \(\Theta(n)\) là độ phức tạp thời gian của
\(\textsc{Partition}(A,lo,hi)\).

Trong trường hợp xấu nhất,
\[
  T(n)=\max_{q=1,2,\ldots,n-1}\{T(q)+T(n-q)\}+\Theta(n).
\]
Có thể chứng minh \(T(n)\le cn^2\), trong đó \(c\) là hằng số, và khi mỗi lần
chia đều làm \(q=1\), có \(T(n)=\Theta(n^2)\). Vì vậy trong trường hợp xấu
nhất, độ phức tạp thời gian của sắp xếp nhanh là \(\Theta(n^2)\).

Chứng minh \(T(n)\le cn^2\) có thể dùng quy nạp toán học. Quá trình tóm tắt:
từ giả thiết quy nạp,
\[
\begin{aligned}
T(n)
&\le \max_{q=1,2,\ldots,n-1}\{cq^2+c(n-q)^2\}+\Theta(n)\\
&= c\cdot \max_{q=1,2,\ldots,n-1}\{q^2+(n-q)^2\}+\Theta(n).
\end{aligned}
\]
Hàm theo \(q\) trong dấu \(\max\) đạt giá trị lớn nhất khi \(q=1\) hoặc
\(q=n-1\), nên
\[
  T(n)\le cn^2-2c(n-1)+\Theta(n)\le cn^2,
\]
trong đó chọn \(c\) thích hợp để triệt tiêu \(-2c(n-1)+\Theta(n)\). Chứng
minh xong.

Trong trường hợp tốt nhất, mỗi lần chia đều chia \(A\) thành hai phần có kích
thước bằng nhau. Khi đó
\[
  T(n)=2T(n/2)+\Theta(n).
\]
Do cây đệ quy hình thành trong quá trình sắp xếp có \(\Theta(\log_2 n)\) tầng,
và chi phí thời gian ở mỗi tầng đều là \(\Theta(n)\), nên
\(T(n)=\Theta(n\log_2 n)\). Phân tích trên bỏ qua trường hợp \(n/2\) không là
số nguyên, nhưng điều đó không ảnh hưởng tới kết quả phân tích.

Ta đã giả sử mọi hoán vị đầu vào xuất hiện với xác suất bằng nhau. Với trường
hợp trung bình,
\[
\begin{aligned}
T(n)
&= \frac{1}{n}\sum_{q=1}^{n-1}(T(q)+T(n-q))+\Theta(n)\\
&= \frac{2}{n}\sum_{q=1}^{n-1}T(q)+\Theta(n).
\end{aligned}
\]
Có thể chứng minh \(T(n)\le an\log_2 n+b\), trong đó \(a,b\) là hằng số, nên
độ phức tạp thời gian trung bình của sắp xếp nhanh là \(O(n\log_2 n)\).

Chứng minh \(T(n)\le an\log_2 n+b\) có thể dùng quy nạp toán học. Quá trình
tóm tắt: từ giả thiết quy nạp,
\[
\begin{aligned}
T(n)
&\le \frac{2}{n}\sum_{q=1}^{n-1}(aq\log_2 q+b)+\Theta(n)\\
&= \frac{2a}{n}\sum_{q=1}^{n-1}q\log_2 q + \frac{2b(n-1)}{n}+\Theta(n).
\end{aligned}
\]
Với hạng thứ nhất,
\[
  \sum_{q=1}^{n-1}q\log_2 q
  = \sum_{q=1}^{\lceil n/2\rceil-1}q\log_2 q
    + \sum_{q=\lceil n/2\rceil}^{n-1}q\log_2 q.
\]
Trong tổng thứ nhất ở vế phải,
\(\log_2 q\le\log_2(n/2)=\log_2 n-1\); trong tổng thứ hai,
\(\log_2 q\le\log_2 n\). Vì vậy,
\[
\begin{aligned}
\sum_{q=1}^{n-1}q\log_2 q
&\le (\log_2 n-1)\sum_{q=1}^{\lceil n/2\rceil-1}q
  + \log_2 n\sum_{q=\lceil n/2\rceil}^{n-1}q\\
&= \log_2 n\sum_{q=1}^{n-1}q
  - \sum_{q=1}^{\lceil n/2\rceil-1}q\\
&\le \frac{n(n-1)}{2}\log_2 n
  - \frac{1}{2}\left(\frac{n}{2}-1\right)\frac{n}{2}\\
&\le \frac{n^2}{2}\log_2 n - \frac{n^2}{8}.
\end{aligned}
\]
Từ đó suy ra
\[
  T(n)\le an\log_2 n+b+\bigl(\Theta(n)+b-an/4\bigr)\le an\log_2 n+b,
\]
trong đó chọn \(a,b\) thích hợp để triệt tiêu
\(\Theta(n)+b-an/4\). Chứng minh xong.

\paragraph{[5]}
Môi trường chạy là Pentium 100MHz, biên dịch bằng BP7.0.

\paragraph{[6]}
Nếu chạy nhiều lần thuật toán ngẫu nhiên hóa có kết quả thực thi xác định,
hiệu suất thời gian có thể rất thấp, nên loại thuật toán ngẫu nhiên hóa này
thường chỉ chạy một lần.

\section*{Tài liệu tham khảo}

\begin{enumerate}
  \item \textit{Thuật toán thực dụng và thiết kế chương trình}, Ngô Văn Hổ và
  cộng sự, Nhà xuất bản Công nghiệp Điện tử.
  \item \textit{Tuyển tập đầy đủ về cấu trúc dữ liệu máy tính và thuật toán
  thực dụng}, Công ty Máy tính Hy Vọng Bắc Kinh, Mưu Nhân chủ biên.
  \item \textit{Toán tổ hợp}, Ngô Văn Hổ và cộng sự, Nhà xuất bản Công nghiệp
  Điện tử.
  \item \textit{Toán tổ hợp}, Lư Khai Trừng, Nhà xuất bản Đại học Thanh Hoa.
  \item \textit{Cấu trúc dữ liệu}, Thi Bá Lạc và cộng sự biên soạn, Nhà xuất
  bản Đại học Phục Đán.
  \item \textit{Dẫn luận thi toán trung học}, Nhà xuất bản Giáo dục Thượng Hải.
\end{enumerate}

\end{document}
